% Options for packages loaded elsewhere
% Options for packages loaded elsewhere
\PassOptionsToPackage{unicode}{hyperref}
\PassOptionsToPackage{hyphens}{url}
\PassOptionsToPackage{dvipsnames,svgnames,x11names}{xcolor}
%
\documentclass[
  letterpaper,
  DIV=11,
  numbers=noendperiod]{scrartcl}
\usepackage{xcolor}
\usepackage{amsmath,amssymb}
\setcounter{secnumdepth}{-\maxdimen} % remove section numbering
\usepackage{iftex}
\ifPDFTeX
  \usepackage[T1]{fontenc}
  \usepackage[utf8]{inputenc}
  \usepackage{textcomp} % provide euro and other symbols
\else % if luatex or xetex
  \usepackage{unicode-math} % this also loads fontspec
  \defaultfontfeatures{Scale=MatchLowercase}
  \defaultfontfeatures[\rmfamily]{Ligatures=TeX,Scale=1}
\fi
\usepackage{lmodern}
\ifPDFTeX\else
  % xetex/luatex font selection
\fi
% Use upquote if available, for straight quotes in verbatim environments
\IfFileExists{upquote.sty}{\usepackage{upquote}}{}
\IfFileExists{microtype.sty}{% use microtype if available
  \usepackage[]{microtype}
  \UseMicrotypeSet[protrusion]{basicmath} % disable protrusion for tt fonts
}{}
\makeatletter
\@ifundefined{KOMAClassName}{% if non-KOMA class
  \IfFileExists{parskip.sty}{%
    \usepackage{parskip}
  }{% else
    \setlength{\parindent}{0pt}
    \setlength{\parskip}{6pt plus 2pt minus 1pt}}
}{% if KOMA class
  \KOMAoptions{parskip=half}}
\makeatother
% Make \paragraph and \subparagraph free-standing
\makeatletter
\ifx\paragraph\undefined\else
  \let\oldparagraph\paragraph
  \renewcommand{\paragraph}{
    \@ifstar
      \xxxParagraphStar
      \xxxParagraphNoStar
  }
  \newcommand{\xxxParagraphStar}[1]{\oldparagraph*{#1}\mbox{}}
  \newcommand{\xxxParagraphNoStar}[1]{\oldparagraph{#1}\mbox{}}
\fi
\ifx\subparagraph\undefined\else
  \let\oldsubparagraph\subparagraph
  \renewcommand{\subparagraph}{
    \@ifstar
      \xxxSubParagraphStar
      \xxxSubParagraphNoStar
  }
  \newcommand{\xxxSubParagraphStar}[1]{\oldsubparagraph*{#1}\mbox{}}
  \newcommand{\xxxSubParagraphNoStar}[1]{\oldsubparagraph{#1}\mbox{}}
\fi
\makeatother

\usepackage{color}
\usepackage{fancyvrb}
\newcommand{\VerbBar}{|}
\newcommand{\VERB}{\Verb[commandchars=\\\{\}]}
\DefineVerbatimEnvironment{Highlighting}{Verbatim}{commandchars=\\\{\}}
% Add ',fontsize=\small' for more characters per line
\usepackage{framed}
\definecolor{shadecolor}{RGB}{241,243,245}
\newenvironment{Shaded}{\begin{snugshade}}{\end{snugshade}}
\newcommand{\AlertTok}[1]{\textcolor[rgb]{0.68,0.00,0.00}{#1}}
\newcommand{\AnnotationTok}[1]{\textcolor[rgb]{0.37,0.37,0.37}{#1}}
\newcommand{\AttributeTok}[1]{\textcolor[rgb]{0.40,0.45,0.13}{#1}}
\newcommand{\BaseNTok}[1]{\textcolor[rgb]{0.68,0.00,0.00}{#1}}
\newcommand{\BuiltInTok}[1]{\textcolor[rgb]{0.00,0.23,0.31}{#1}}
\newcommand{\CharTok}[1]{\textcolor[rgb]{0.13,0.47,0.30}{#1}}
\newcommand{\CommentTok}[1]{\textcolor[rgb]{0.37,0.37,0.37}{#1}}
\newcommand{\CommentVarTok}[1]{\textcolor[rgb]{0.37,0.37,0.37}{\textit{#1}}}
\newcommand{\ConstantTok}[1]{\textcolor[rgb]{0.56,0.35,0.01}{#1}}
\newcommand{\ControlFlowTok}[1]{\textcolor[rgb]{0.00,0.23,0.31}{\textbf{#1}}}
\newcommand{\DataTypeTok}[1]{\textcolor[rgb]{0.68,0.00,0.00}{#1}}
\newcommand{\DecValTok}[1]{\textcolor[rgb]{0.68,0.00,0.00}{#1}}
\newcommand{\DocumentationTok}[1]{\textcolor[rgb]{0.37,0.37,0.37}{\textit{#1}}}
\newcommand{\ErrorTok}[1]{\textcolor[rgb]{0.68,0.00,0.00}{#1}}
\newcommand{\ExtensionTok}[1]{\textcolor[rgb]{0.00,0.23,0.31}{#1}}
\newcommand{\FloatTok}[1]{\textcolor[rgb]{0.68,0.00,0.00}{#1}}
\newcommand{\FunctionTok}[1]{\textcolor[rgb]{0.28,0.35,0.67}{#1}}
\newcommand{\ImportTok}[1]{\textcolor[rgb]{0.00,0.46,0.62}{#1}}
\newcommand{\InformationTok}[1]{\textcolor[rgb]{0.37,0.37,0.37}{#1}}
\newcommand{\KeywordTok}[1]{\textcolor[rgb]{0.00,0.23,0.31}{\textbf{#1}}}
\newcommand{\NormalTok}[1]{\textcolor[rgb]{0.00,0.23,0.31}{#1}}
\newcommand{\OperatorTok}[1]{\textcolor[rgb]{0.37,0.37,0.37}{#1}}
\newcommand{\OtherTok}[1]{\textcolor[rgb]{0.00,0.23,0.31}{#1}}
\newcommand{\PreprocessorTok}[1]{\textcolor[rgb]{0.68,0.00,0.00}{#1}}
\newcommand{\RegionMarkerTok}[1]{\textcolor[rgb]{0.00,0.23,0.31}{#1}}
\newcommand{\SpecialCharTok}[1]{\textcolor[rgb]{0.37,0.37,0.37}{#1}}
\newcommand{\SpecialStringTok}[1]{\textcolor[rgb]{0.13,0.47,0.30}{#1}}
\newcommand{\StringTok}[1]{\textcolor[rgb]{0.13,0.47,0.30}{#1}}
\newcommand{\VariableTok}[1]{\textcolor[rgb]{0.07,0.07,0.07}{#1}}
\newcommand{\VerbatimStringTok}[1]{\textcolor[rgb]{0.13,0.47,0.30}{#1}}
\newcommand{\WarningTok}[1]{\textcolor[rgb]{0.37,0.37,0.37}{\textit{#1}}}

\usepackage{longtable,booktabs,array}
\newcounter{none} % for unnumbered tables
\usepackage{calc} % for calculating minipage widths
% Correct order of tables after \paragraph or \subparagraph
\usepackage{etoolbox}
\makeatletter
\patchcmd\longtable{\par}{\if@noskipsec\mbox{}\fi\par}{}{}
\makeatother
% Allow footnotes in longtable head/foot
\IfFileExists{footnotehyper.sty}{\usepackage{footnotehyper}}{\usepackage{footnote}}
\makesavenoteenv{longtable}
\usepackage{graphicx}
\makeatletter
\newsavebox\pandoc@box
\newcommand*\pandocbounded[1]{% scales image to fit in text height/width
  \sbox\pandoc@box{#1}%
  \Gscale@div\@tempa{\textheight}{\dimexpr\ht\pandoc@box+\dp\pandoc@box\relax}%
  \Gscale@div\@tempb{\linewidth}{\wd\pandoc@box}%
  \ifdim\@tempb\p@<\@tempa\p@\let\@tempa\@tempb\fi% select the smaller of both
  \ifdim\@tempa\p@<\p@\scalebox{\@tempa}{\usebox\pandoc@box}%
  \else\usebox{\pandoc@box}%
  \fi%
}
% Set default figure placement to htbp
\def\fps@figure{htbp}
\makeatother





\setlength{\emergencystretch}{3em} % prevent overfull lines

\providecommand{\tightlist}{%
  \setlength{\itemsep}{0pt}\setlength{\parskip}{0pt}}



 


\KOMAoption{captions}{tableheading}
\makeatletter
\@ifpackageloaded{tcolorbox}{}{\usepackage[skins,breakable]{tcolorbox}}
\@ifpackageloaded{fontawesome5}{}{\usepackage{fontawesome5}}
\definecolor{quarto-callout-color}{HTML}{909090}
\definecolor{quarto-callout-note-color}{HTML}{0758E5}
\definecolor{quarto-callout-important-color}{HTML}{CC1914}
\definecolor{quarto-callout-warning-color}{HTML}{EB9113}
\definecolor{quarto-callout-tip-color}{HTML}{00A047}
\definecolor{quarto-callout-caution-color}{HTML}{FC5300}
\definecolor{quarto-callout-color-frame}{HTML}{acacac}
\definecolor{quarto-callout-note-color-frame}{HTML}{4582ec}
\definecolor{quarto-callout-important-color-frame}{HTML}{d9534f}
\definecolor{quarto-callout-warning-color-frame}{HTML}{f0ad4e}
\definecolor{quarto-callout-tip-color-frame}{HTML}{02b875}
\definecolor{quarto-callout-caution-color-frame}{HTML}{fd7e14}
\makeatother
\makeatletter
\@ifpackageloaded{caption}{}{\usepackage{caption}}
\AtBeginDocument{%
\ifdefined\contentsname
  \renewcommand*\contentsname{Table of contents}
\else
  \newcommand\contentsname{Table of contents}
\fi
\ifdefined\listfigurename
  \renewcommand*\listfigurename{List of Figures}
\else
  \newcommand\listfigurename{List of Figures}
\fi
\ifdefined\listtablename
  \renewcommand*\listtablename{List of Tables}
\else
  \newcommand\listtablename{List of Tables}
\fi
\ifdefined\figurename
  \renewcommand*\figurename{Figure}
\else
  \newcommand\figurename{Figure}
\fi
\ifdefined\tablename
  \renewcommand*\tablename{Table}
\else
  \newcommand\tablename{Table}
\fi
}
\@ifpackageloaded{float}{}{\usepackage{float}}
\floatstyle{ruled}
\@ifundefined{c@chapter}{\newfloat{codelisting}{h}{lop}}{\newfloat{codelisting}{h}{lop}[chapter]}
\floatname{codelisting}{Listing}
\newcommand*\listoflistings{\listof{codelisting}{List of Listings}}
\makeatother
\makeatletter
\makeatother
\makeatletter
\@ifpackageloaded{caption}{}{\usepackage{caption}}
\@ifpackageloaded{subcaption}{}{\usepackage{subcaption}}
\makeatother
\usepackage{bookmark}
\IfFileExists{xurl.sty}{\usepackage{xurl}}{} % add URL line breaks if available
\urlstyle{same}
\hypersetup{
  pdftitle={Introducing Time Series Analysis},
  colorlinks=true,
  linkcolor={blue},
  filecolor={Maroon},
  citecolor={Blue},
  urlcolor={Blue},
  pdfcreator={LaTeX via pandoc}}


\title{Introducing Time Series Analysis}
\usepackage{etoolbox}
\makeatletter
\providecommand{\subtitle}[1]{% add subtitle to \maketitle
  \apptocmd{\@title}{\par {\large #1 \par}}{}{}
}
\makeatother
\subtitle{Topic 1}
\author{}
\date{}
\begin{document}
\maketitle


\begin{center}\rule{0.5\linewidth}{0.5pt}\end{center}

{Goals}

\begin{itemize}
\tightlist
\item
  To review basic, but crucial, statistical concepts
\item
  To introduce what is time series analysis
\item
  To motivate why we need special techniques for analysing time series
\item
  To start exercising our analytical skills with algebra that steadily
  increases in complexity
\end{itemize}

\subsection{Background knowledge}\label{background-knowledge}

\subsubsection{Statistics}\label{statistics}

It is assumed that you have previously studied basic probability theory
and statistical inference and have a sound understanding of the
following:

\begin{itemize}
\tightlist
\item
  Distinction between a population and a sample
\item
  Descriptive statistics --- sample mean, median, quartiles, variance
\item
  Random variables and stochastic processes
\item
  CDFs, PMFs, and PDFs of random variables
\item
  Expected values and linearity of the \(\mathsf{E}(\cdot)\) operator
\item
  Population variance, covariance, and correlation
\item
  Sampling theory, estimation, unbiasedness, and efficiency
\item
  Mean square error and Markov/Chebyshev inequalities
\item
  Normal distribution
\item
  Hypothesis testing --- Type I and II errors, significance level,
  power, \(p\)-value
\item
  Confidence intervals
\item
  Convergence in probability, LLN, \(p\)lim operator, and consistency
\item
  Convergence in distribution, CLT, and limiting distributions
\end{itemize}

\subsubsection{Econometrics}\label{econometrics}

It is assumed that you have previously studied basic econometrics and
have familiarity with:

\begin{itemize}
\tightlist
\item
  Multiple linear regression models --- classical and generalised
  settings
\item
  Estimation via OLS, MME, and MLE
\item
  Asymptotic evaluation of estimators --- consistency, asymptotic
  efficiency, asymptotic normality
\end{itemize}

\begin{tcolorbox}[enhanced jigsaw, arc=.35mm, bottomrule=.15mm, bottomtitle=1mm, breakable, colback=white, colbacktitle=quarto-callout-important-color!10!white, colframe=quarto-callout-important-color-frame, coltitle=black, left=2mm, leftrule=.75mm, opacityback=0, opacitybacktitle=0.6, rightrule=.15mm, title=\textcolor{quarto-callout-important-color}{\faExclamation}\hspace{0.5em}{Important}, titlerule=0mm, toprule=.15mm, toptitle=1mm]

\begin{itemize}
\tightlist
\item
  If any of these topics are unclear to you, please seek my help in
  office hours as early in the term as possible.
\item
  Alternatively, review any good introductory statistics/econometrics
  book of your choice. I can suggest \emph{Larsen and Marx} and/or
  \emph{Jeffrey Wooldridge} for self-study.
\end{itemize}

\end{tcolorbox}

\subsection{Basic concepts}\label{basic-concepts}

{What is a random variable?}

\begin{itemize}
\item
  A random variable is a \textbf{function}
  \(Y : \Omega \mapsto \mathbb{R}\) that maps the outcomes of an
  experiment to a real number.
\item
  Random variables are denoted by capital italicised Roman letters such
  as \(Y = Y(\omega)\), where \(\omega \in \Omega\).
\item
  For example, suppose \(H\) and \(T\) are the two possible outcomes of
  a coin-tossing experiment. The statement
  \begin{align} Y = 1 \text{ if }H\text{ occurs and }Y = 0\text{ if }T\text{ occurs} \end{align}
  defines a random variable.
\end{itemize}

A real-valued function of the sample space.

{What is a stochastic process?}

\begin{itemize}
\item
  Consider a \textbf{collection} of random variables denoted by the set
  \[\{Y_i : i \in \mathbb{Z}\} \quad \text{or} \quad \{Y_i\}_{i \in \mathbb{Z}},\]
  commonly referred to as a stochastic process. The indexation \(i\)
  represents an arbitrary dimension along which the random variables are
  ordered within the collection.
\item
  Given data \(\{y_1, \ldots, y_N\}\), the standard approach models the
  dataset as a finite-length \(N\) realisation of this stochastic
  process.
\item
  The relationship among the random variables is one of its critical
  defining features --- for instance, \(\{y_1, \ldots, y_N\}\) is a
  random sample if the random variables are iid.
\end{itemize}

A collection of random variables.

{What is a time series?}

\begin{itemize}
\item
  Consider what happens when we switch notation from \(i\) to \(t\),
  characterising our dataset as \(\{y_t\}_{t=1}^{N}\) instead of
  \(\{y_i\}_{i=1}^{N}\).
\item
  We now think of the subscript as denoting a point in time. The
  sampling dimension is now \textbf{temporal}.
\item
  Algebraically this switch looks inconsequential. Conceptually, the
  change is monumental.
\item
  If we think of the data as a collection of observations over time,
  three issues immediately arise that cannot be ignored from a modelling
  perspective.
\end{itemize}

A stochastic process indexed over time.

{What key issues emerge when considering time series?}

\textbf{Issue 1 --- Ordering.} There exists a natural fixed ordering of
the data. Re-shuffling, for whatever statistical reason, would not be
advisable.

\textbf{Issue 2 --- Stability.} There exists a clear logical requirement
to characterise stability of the data-generating process over time.
Otherwise our sample is an amalgamation of data from distinct
probabilistic regimes.

\begin{tcolorbox}[enhanced jigsaw, arc=.35mm, bottomrule=.15mm, bottomtitle=1mm, breakable, colback=white, colbacktitle=quarto-callout-note-color!10!white, colframe=quarto-callout-note-color-frame, coltitle=black, left=2mm, leftrule=.75mm, opacityback=0, opacitybacktitle=0.6, rightrule=.15mm, title=\textcolor{quarto-callout-note-color}{\faInfo}\hspace{0.5em}{Watch: Illustration}, titlerule=0mm, toprule=.15mm, toptitle=1mm]

\textbf{Estimation under Instability}

\end{tcolorbox}

\textbf{Issue 3 --- Dependence.} There exists an overriding need to
manage the presence of dependence. One may be willing to assume
independence over the cross-section, but one could never justify
independence over time. Dependence is the defining feature of the field
of time series analysis.

Issues 2 and 3 may seem abstract at present, but their meaning will
become clear when we discuss stationarity and ergodicity later.

{What is the defining feature of (stable) time series?}

\begin{itemize}
\item
  Consider again an earlier statement: dependence is the defining
  feature of the field of time series analysis. What does this statement
  mean?
\item
  It means that we will distinguish one time series model from another
  by reference precisely to their dependence characteristics.
\item
  In particular, our primary tool will be with so-called
  autocovariances. These functions (of lag length) capture the linear
  dependence characteristics between the random variables comprising a
  time series.
\end{itemize}

{We characterise time series via their dependence characteristics. In
this course, as in most introductory courses, we will focus solely on
linear dependence (autocovariances).}

{Illustrative examples of (stable) time series}

\begin{itemize}
\item
  An \(iid\) process \(\{Y_t\}_{t \in \mathbb{Z}}\) has no linear or
  non-linear dependence on the past --- since the random variables are
  independent, the lag \(h\) autocovariance,
  \(\gamma_Y(h) := \text{Cov}(Y_t, Y_{t+h})\), satisfies
  \(\gamma_Y(h) = 0\) for any \(h \neq 0\). We will later refer to an
  \(iid\) process as \textbf{strong white noise}.
\item
  By contrast, consider a process where
  \(Y_t = \varepsilon_t + 0.5\varepsilon_{t-1}\) and
  \(\{\varepsilon_t\}\) is strong white noise. We have
  \begin{align} \gamma_Y(-1) &= \text{Cov}(Y_t, Y_{t-1})\\ &= \text{Cov}(\varepsilon_t + 0.5\varepsilon_{t-1},\, \varepsilon_{t-1} + 0.5\varepsilon_{t-2})\\ &= 0.5\,\text{Var}(\varepsilon_{t-1}) \text{ since }\{\varepsilon_t\}\text{ is $iid$}\\ &\neq 0.\end{align}
  Notice that \(\gamma_Y(1) = \gamma_Y(-1) \neq 0\) and
  \(\gamma_Y(h) = 0\) for all \(|h| \geq 2\). We will later refer to
  this as a \textbf{moving average process of order 1}.
\end{itemize}

These are examples of processes with and without serial correlation,
respectively.

\begin{tcolorbox}[enhanced jigsaw, arc=.35mm, bottomrule=.15mm, bottomtitle=1mm, breakable, colback=white, colbacktitle=quarto-callout-note-color!10!white, colframe=quarto-callout-note-color-frame, coltitle=black, left=2mm, leftrule=.75mm, opacityback=0, opacitybacktitle=0.6, rightrule=.15mm, title=\textcolor{quarto-callout-note-color}{\faInfo}\hspace{0.5em}{Watch: Live Simulation}, titlerule=0mm, toprule=.15mm, toptitle=1mm]

\begin{Shaded}
\begin{Highlighting}[]
\NormalTok{\#| \textquotesingle{}!! shinylive warning !!\textquotesingle{}: |}
\NormalTok{\#|   shinylive does not work in self{-}contained HTML documents.}
\NormalTok{\#|   Please set \textasciigrave{}embed{-}resources: false\textasciigrave{} in your metadata.}
\NormalTok{\#| standalone: true}
\NormalTok{\#| viewerHeight: 560}

\NormalTok{library(shiny)}

\NormalTok{NAVY   \textless{}{-} "\#0D1B2A"}
\NormalTok{AMBER  \textless{}{-} "\#C8762A"}
\NormalTok{BG     \textless{}{-} "\#FAFAF8"}
\NormalTok{BG2    \textless{}{-} "\#F5F0E8"}
\NormalTok{INK    \textless{}{-} "\#2C2420"}
\NormalTok{GRAPH  \textless{}{-} "\#8A7A68"}
\NormalTok{STONE  \textless{}{-} "\#B0A494"}
\NormalTok{RULE   \textless{}{-} "\#E0D8CC"}
\NormalTok{BORDER \textless{}{-} "\#D8CEBC"}

\NormalTok{ui \textless{}{-} fluidPage(}

\NormalTok{  tags$head(tags$style(HTML(paste0("}
\NormalTok{    body, .shiny{-}frame \{ background:", BG, "; color:", INK, "; \}}

\NormalTok{    .ctrl{-}panel \{}
\NormalTok{      background:", BG2, ";}
\NormalTok{      border: 0.5px solid ", BORDER, ";}
\NormalTok{      border{-}top: 3px solid ", AMBER, ";}
\NormalTok{      border{-}radius: 4px;}
\NormalTok{      padding: 1.1rem 1.2rem;}
\NormalTok{      margin{-}top: 0.5rem;}
\NormalTok{    \}}
\NormalTok{    .ctrl{-}label \{}
\NormalTok{      font{-}family: \textquotesingle{}Courier New\textquotesingle{}, monospace;}
\NormalTok{      font{-}size: 0.85rem;}
\NormalTok{      letter{-}spacing: 0.15em;}
\NormalTok{      color: ", STONE, ";}
\NormalTok{      margin{-}bottom: 0.15rem;}
\NormalTok{      display: block;}
\NormalTok{    \}}
\NormalTok{    .gamma{-}val \{}
\NormalTok{      font{-}family: \textquotesingle{}Courier New\textquotesingle{}, monospace;}
\NormalTok{      font{-}size: 0.82rem;}
\NormalTok{      color: ", AMBER, ";}
\NormalTok{      margin{-}top: {-}0.3rem;}
\NormalTok{      margin{-}bottom: 0.6rem;}
\NormalTok{    \}}
\NormalTok{    .ctrl{-}divider \{}
\NormalTok{      border: none;}
\NormalTok{      border{-}top: 0.5px solid ", RULE, ";}
\NormalTok{      margin: 0.75rem 0;}
\NormalTok{    \}}
\NormalTok{    .irs{-}{-}shiny .irs{-}bar \{ background:", NAVY, "; border{-}color:", NAVY, "; \}}
\NormalTok{    .irs{-}{-}shiny .irs{-}handle \{ border{-}color:", AMBER, "; \}}
\NormalTok{    .irs{-}{-}shiny .irs{-}single \{ background:", NAVY, "; \}}
\NormalTok{    .irs{-}{-}shiny .irs{-}from,}
\NormalTok{    .irs{-}{-}shiny .irs{-}to   \{ background:", NAVY, "; \}}
\NormalTok{    .irs{-}{-}shiny .irs{-}grid{-}pol \{ background:", RULE, "; \}}
\NormalTok{    .irs{-}{-}shiny .irs{-}line \{ background:", RULE, "; border{-}color:", RULE, "; \}}
\NormalTok{    label.control{-}label \{ display: none; \}}
\NormalTok{  ")))),}

\NormalTok{  fluidRow(}
\NormalTok{    column(12,}
\NormalTok{      plotOutput("upper\_plot", height = "230px"),}
\NormalTok{      div(style = "height:6px"),}
\NormalTok{      plotOutput("lower\_plot", height = "230px")}
\NormalTok{    )}
\NormalTok{  ),}
\NormalTok{  fluidRow(}
\NormalTok{    column(12,}
\NormalTok{      div(class = "ctrl{-}panel",}
\NormalTok{        span(class = "ctrl{-}label", style = "font{-}size: 1.5rem;", "Coefficient"),}
\NormalTok{        div(style = "width: 100\%;",}
\NormalTok{          sliderInput("theta",}
\NormalTok{                      label = NULL, min = {-}1, max = 1, value = 0.5,}
\NormalTok{                      step = 0.05, ticks = FALSE, width = "100\%")}
\NormalTok{        )}
\NormalTok{      )}
\NormalTok{    )}
\NormalTok{  )}
\NormalTok{)}

\NormalTok{server \textless{}{-} function(input, output, session) \{}

\NormalTok{  series \textless{}{-} reactive(\{}
\NormalTok{    set.seed(42L)}
\NormalTok{    n     \textless{}{-} 100L}
\NormalTok{    theta \textless{}{-} input$theta}
\NormalTok{    eps   \textless{}{-} rnorm(n + 1L)}
\NormalTok{    list(}
\NormalTok{      iid = eps[{-}1L],}
\NormalTok{      ma1 = eps[{-}1L] + theta * eps[{-}(n + 1L)],}
\NormalTok{      t   = seq\_len(n)}
\NormalTok{    )}
\NormalTok{  \})}

\NormalTok{  output$gamma\_display \textless{}{-} renderText(\{}
\NormalTok{    paste0("γ(1) = ", format(input$theta, nsmall = 2))}
\NormalTok{  \})}

\NormalTok{  base\_par \textless{}{-} function() \{}
\NormalTok{    par(}
\NormalTok{      bg        = BG,}
\NormalTok{      col.axis  = GRAPH,}
\NormalTok{      col.lab   = INK,}
\NormalTok{      col.main  = NAVY,}
\NormalTok{      family    = "serif",}
\NormalTok{      mar       = c(3, 3.8, 2.2, 0.8),}
\NormalTok{      mgp       = c(2.2, 0.5, 0),}
\NormalTok{      tcl       = {-}0.25,}
\NormalTok{      cex.axis  = 0.75,}
\NormalTok{      cex.lab   = 0.82,}
\NormalTok{      cex.main  = 0.88,}
\NormalTok{      font.main = 1}
\NormalTok{    )}
\NormalTok{  \}}

\NormalTok{  ylims \textless{}{-} reactive(\{}
\NormalTok{    s   \textless{}{-} series()}
\NormalTok{    rng \textless{}{-} range(c(s$iid, s$ma1))}
\NormalTok{    pad \textless{}{-} diff(rng) * 0.06}
\NormalTok{    c(rng[1] {-} pad, rng[2] + pad)}
\NormalTok{  \})}

\NormalTok{  output$upper\_plot \textless{}{-} renderPlot(\{}
\NormalTok{    s \textless{}{-} series()}
\NormalTok{    base\_par()}
\NormalTok{    plot(s$t, s$iid,}
\NormalTok{         type = "l", lwd = 1.0, col = NAVY,}
\NormalTok{         ylim = ylims(),}
\NormalTok{         main = expression(paste("iid  ", N(0,1), "  noise")),}
\NormalTok{         xlab = "", ylab = expression(Y[t]),}
\NormalTok{         bty = "l", xaxt = "n")}
\NormalTok{    abline(h = 0, col = RULE, lwd = 0.9)}
\NormalTok{  \}, bg = BG)}

\NormalTok{  output$lower\_plot \textless{}{-} renderPlot(\{}
\NormalTok{    s     \textless{}{-} series()}
\NormalTok{    theta \textless{}{-} input$theta}
\NormalTok{    base\_par()}
\NormalTok{    plot(s$t, s$ma1,}
\NormalTok{         type = "l", lwd = 1.2, col = AMBER,}
\NormalTok{         ylim = ylims(),}
\NormalTok{         main = bquote(}
\NormalTok{           MA(1): \textasciitilde{} Y[t] == epsilon[t] + .(theta) \textasciitilde{} epsilon[t{-}1] \textasciitilde{}}
\NormalTok{           "    " \textasciitilde{} gamma(1) == .(theta)}
\NormalTok{         ),}
\NormalTok{         xlab = expression(t), ylab = expression(Y[t]),}
\NormalTok{         bty = "l")}
\NormalTok{    abline(h = 0, col = RULE, lwd = 0.9)}
\NormalTok{  \}, bg = BG)}
\NormalTok{\}}

\NormalTok{shinyApp(ui, server)}
\end{Highlighting}
\end{Shaded}

\end{tcolorbox}

Students sometimes ask why I present an \(iid\) process as a time series
example after emphasising that time series are characterised by
dependence. The answer: an \(iid\) process indexed by time is a special
case of a time series --- one where the dependence happens to be zero.
Whether you call it a ``genuine'' time series is a matter of
perspective.

{What is time series analysis?}

\begin{itemize}
\item
  We will soon see (see ``\hyperref[a-stark-message]{A stark message}''
  below) that the analysis of time series data can lead to new and
  unique problems in statistical modelling and inference.
\item
  Specifically, the obvious possibility of dependence introduced by the
  sampling of adjacent points in time can severely restrict the
  applicability of the many conventional statistical methods
  traditionally dependent on the assumption that these adjacent
  observations are independent and identically distributed.
\end{itemize}

The systematic approach by which one goes about answering econometric
questions posed in the presence of time-ordered data exhibiting
inter-temporal dependence is commonly referred to as time series
analysis.

Two approaches to time series analysis exist. The \textbf{time domain}
approach studies lagged relationships --- e.g.~how today's value affects
tomorrow's. The \textbf{frequency domain} approach studies cycles by, in
effect, regressing the series on sinusoids at different frequencies.
This module focuses exclusively on the time domain approach.

\begin{tcolorbox}[enhanced jigsaw, arc=.35mm, bottomrule=.15mm, bottomtitle=1mm, breakable, colback=white, colbacktitle=quarto-callout-note-color!10!white, colframe=quarto-callout-note-color-frame, coltitle=black, left=2mm, leftrule=.75mm, opacityback=0, opacitybacktitle=0.6, rightrule=.15mm, title=\textcolor{quarto-callout-note-color}{\faInfo}\hspace{0.5em}{Watch: Illustration}, titlerule=0mm, toprule=.15mm, toptitle=1mm]

\textbf{Time vs.~Frequency Domains}

\end{tcolorbox}

\begin{center}\rule{0.5\linewidth}{0.5pt}\end{center}

\subsection{Variance, covariance, and
correlation}\label{variance-covariance-and-correlation}

Covariances --- and their time-series counterpart, autocovariances ---
are central to everything that follows. In this section we recap the key
concepts summarising the second-order structure of the probabilistic
behaviour of random variables.

\subsubsection{Definitions}\label{definitions}

Consider univariate random variables \(Y_1\) and \(Y_2\). The \(r\)-th
\textbf{raw moment} of \(Y_1\) is \(\mathsf{E}(Y_1^r)\) and the \(r\)-th
\textbf{central moment} is
\(\mathsf{E}\!\left[(Y_1 - \mathsf{E}(Y_1))^r\right]\), for
\(r \in \mathbb{N}\).

\[\text{Var}(Y_1) := \mathsf{E}\!\left[(Y_1 - \mathsf{E}(Y_1))^2\right] = \mathsf{E}(Y_1^2) - (\mathsf{E}(Y_1))^2\]

The variance is the second central moment of \(Y_1\) --- a measure of
dispersion.

\[\text{Cov}(Y_1, Y_2) := \mathsf{E}\!\left[(Y_1 - \mathsf{E}(Y_1))(Y_2 - \mathsf{E}(Y_2))\right] = \mathsf{E}(Y_1 Y_2) - \mathsf{E}(Y_1)\mathsf{E}(Y_2)\]

Note that \(\text{Var}(Y_1) = \text{Cov}(Y_1, Y_1)\) trivially ---
variance is a special case of covariance where both indices are equal.

\[\text{Corr}(Y_1, Y_2) := \frac{\text{Cov}(Y_1, Y_2)}{\sqrt{\text{Var}(Y_1)\,\text{Var}(Y_2)}}\]

Correlation is a scale-free measure of the strength of linear
association, and \(|\text{Corr}(Y_1, Y_2)| \leq 1\) by the
Cauchy-Schwarz inequality.

\begin{center}\rule{0.5\linewidth}{0.5pt}\end{center}

{What is the variance of a linear combination?}

audio/intro.mp3

Consider random variables \(Y_1\) and \(Y_2\) and constants \(a_1\) and
\(a_2\). The familiar variance rule states

\[\text{Var}(a_1 Y_1 + a_2 Y_2) = a_1^2\,\text{Var}(Y_1) + 2a_1 a_2\,\text{Cov}(Y_1, Y_2) + a_2^2\,\text{Var}(Y_2).\]
But if you allow me to rewrite the above as:

\[\text{Var}(a_1 Y_1 + a_2 Y_2) = a_1 a_1 \,\text{Cov}(Y_1,Y_1) + a_1 a_2\,\text{Cov}(Y_1, Y_2) + a_2 a_1\,\text{Cov}(Y_2, Y_1) +  a_2a_2\,\text{Cov}(Y_2,Y_2),\]
then you can see that we can also use summation notation to express the
result as:

\[\text{Var}(a_1 Y_1 + a_2 Y_2) = \sum_{t=1}^{2}\sum_{s=1}^{2} a_t a_s\,\text{Cov}(Y_t, Y_s).\]

Indeed, it should come as no surprise then that

\[\text{Var}\!\left(\sum_{t=1}^{T} a_t Y_t\right) = \sum_{t=1}^{T}\sum_{s=1}^{T} a_t a_s\,\text{Cov}(Y_t, Y_s).\]

which is an expression we will see repeatedly.

Never omit the covariance terms!

\begin{tcolorbox}[enhanced jigsaw, arc=.35mm, bottomrule=.15mm, bottomtitle=1mm, breakable, colback=white, colbacktitle=quarto-callout-note-color!10!white, colframe=quarto-callout-note-color-frame, coltitle=black, left=2mm, leftrule=.75mm, opacityback=0, opacitybacktitle=0.6, rightrule=.15mm, title=\textcolor{quarto-callout-note-color}{\faInfo}\hspace{0.5em}{Intuition}, titlerule=0mm, toprule=.15mm, toptitle=1mm]

\begin{itemize}
\item
  Students don't always find variance/covariance manipulations
  intuitive. However, the same students would never make the mistake of
  omitting cross-terms in a squaring operation! The sum of two terms
  squared must yield a sum of four terms. Three terms must yield nine.
  Let's look at an example:
  \[(a + b + c)^2 = \begin{align}&a^2 + ab + ac \\+ &ba + b^2 + bc \\ + &ca + cb + c^2\end{align}\]
  This is nine terms, right? You would never say
  \[(a+b+c)^2 = a^2 + b^2 + c^2.\quad \text{❌}\] It just looks wrong!
\item
  But students make this mistake all the time with variances. They say
  \[\mathsf{Var}(X+Y) = \mathsf{Var}(X) + \mathsf{Var}(Y).\quad \text{❌}\]

  Don't forget that a variance is a squared concept (see algebra below).
  You cannot just assume that the cross-terms are zero unless you are
  told or you know explicitly that they are zero.
\end{itemize}

\end{tcolorbox}

\begin{tcolorbox}[enhanced jigsaw, arc=.35mm, bottomrule=.15mm, bottomtitle=1mm, breakable, colback=white, colbacktitle=quarto-callout-note-color!10!white, colframe=quarto-callout-note-color-frame, coltitle=black, left=2mm, leftrule=.75mm, opacityback=0, opacitybacktitle=0.6, rightrule=.15mm, title=\textcolor{quarto-callout-note-color}{\faInfo}\hspace{0.5em}{Algebra}, titlerule=0mm, toprule=.15mm, toptitle=1mm]

\begin{itemize}
\tightlist
\item
  Consider \(\mathsf{Var}(Z) := \mathsf{E}(Z-\mathsf{E}(Z))^2\). Now
  suppose \(Z = X+Y\). Then,
  \begin{align}\mathsf{Var}(Z)=\mathsf{Var}(X+Y)&=\mathsf{E}\left((X+Y) - \mathsf{E}(X+Y)\right)^2 \\ &= \mathsf{E}(X - \mathsf{E}(X) + Y - \mathsf{E}(Y))^2\\&=\mathsf{E}\left(\tilde{X} + \tilde{Y}\right)^2, \text{ where }\tilde{X} := X- \mathsf{E}(X)\text{ and }\tilde{Y} := Y - \mathsf{E}(Y)\\&=\mathsf{E}\left(\tilde{X}^2 + \tilde{Y}^2 + 2\tilde{X}\tilde{Y}\right) \\&=\mathsf{E}\left(\tilde{X}^2\right) + \mathsf{E}\left(\tilde{Y}^2\right) + 2\mathsf{E}\left(\tilde{X}\tilde{Y}\right)\\&=\mathsf{E}(X-\mathsf{E}(X))^2 +\mathsf{E}(Y-\mathsf{E}(Y))^2 + 2 \mathsf{E}((X-\mathsf{E}(X))(Y-\mathsf{E}(Y)))\\&=\mathsf{Var}(X)+\mathsf{Var}(Y) + 2\mathsf{Cov}(X,Y).\end{align}
\item
  Never omit the covariance terms!
\end{itemize}

\end{tcolorbox}

We are now ready for the final section of Topic 1.

\begin{center}\rule{0.5\linewidth}{0.5pt}\end{center}

\subsection{A stark message}\label{a-stark-message}

{What is the sample mean?}

Recall the humble yet ubiquitous sample mean \(\bar{Y}_N\):

\[\bar{Y}_N = \frac{1}{N}\sum_{i=1}^{N} Y_i\]

The subscript \(N\) appears explicitly to emphasise the dependence of
the sample mean on sample size.

Why do we turn so often to this statistic? The answer is
\textbf{Chebyshev's LLN}: \(\bar{Y}_N\) converges in probability to
\(\mathsf{E}(Y)\) as \(N \to \infty\). Under suitable conditions, a
sample moment is a consistent estimator of its population counterpart.

The sample mean appears embedded throughout econometric estimation. For
OLS in \(Y_t = \alpha + \beta X_t + \varepsilon_t\):

\[\hat{\alpha}_{\scriptscriptstyle OLS} = \bar{Y}_T - \hat{\beta}_{\scriptscriptstyle OLS}\bar{X}_T, \qquad \hat{\beta}_{\scriptscriptstyle OLS} = \frac{\frac{1}{T}\sum_{t=1}^{T}(X_t - \bar{X}_T)(Y_t - \bar{Y}_T)}{\frac{1}{T}\sum_{t=1}^{T}(X_t - \bar{X}_T)^2}\]

\((1/N)\sum Y_i\); the most widely-used statistic in econometrics.

{Why is the sample mean ubiquitous? Chebyshev's LLN}

The generic \(F_Y\) notation denotes deliberate agnosticism about the
distribution --- the result holds for any distribution with finite
variance.

\textbf{Theorem --- Chebyshev's LLN.}
\[\text{Suppose } Y_t \overset{\scriptscriptstyle\text{iid}}{\sim} F_Y(\mu_Y, \sigma_Y^2) \text{ for } t = 1,\ldots,T \text{ where } 0 < \sigma_Y^2 < \infty. \text{ Then } \operatorname*{plim}_{T\to\infty}\bar{Y}_T = \mu_Y.\]

Estimator \(\hat{\theta}_T\) \textbf{converges in probability} to
\(\theta\) if and only if for any \(e > 0\):

\[\lim_{T\to\infty} \Pr\!\left(|\hat{\theta}_T - \theta| > e\right) = 0.\]

We write \(\hat{\theta}_T \overset{p}{\to} \theta\) and say
\(\hat{\theta}_T\) is \textbf{consistent} for \(\theta\). Convergence in
mean-square implies convergence in probability via Chebyshev's
inequality:

\[\lim_{T\to\infty} \text{MSE}(\hat{\theta}_T) = \lim_{T\to\infty}\left[\text{Var}(\hat{\theta}_T) + \text{Bias}^2(\hat{\theta}_T)\right] = 0 \implies \hat{\theta}_T \overset{p}{\to} \theta.\]

If \(\text{MSE}(\hat\theta_T) \to 0\) then
\(\hat\theta_T \overset{p}{\to} \theta\); the sample mean achieves this
because \(\text{Var}(\bar{Y}_T) = \sigma^2_Y / T \to 0\).

{What is the mechanism behind Chebyshev's LLN?}

Consider \(\bar{Y}_T = (1/T)\sum_{t=1}^{T} Y_t\). By linearity of
expectations:

\[\mathsf{E}(\bar{Y}_T) = \frac{1}{T}\sum_{t=1}^{T}\mathsf{E}(Y_t) = \frac{T\mu_Y}{T} = \mu_Y,\]

so \(\bar{Y}_T\) is unbiased for \(\mu_Y\). For the variance, using
independence then identical distributions:

\[\text{Var}(\bar{Y}_T) = \frac{1}{T^2}\sum_{t=1}^{T}\sum_{s=1}^{T}\text{Cov}(Y_t, Y_s) = \frac{1}{T^2}\sum_{t=1}^{T}\text{Cov}(Y_t, Y_t) = \frac{\sigma_Y^2}{T}.\]

It follows that
\(\lim_{T\to\infty}\text{MSE}(\bar{Y}_T) = \lim_{T\to\infty}\sigma_Y^2/T = 0\).
By Chebyshev's inequality, for any \(e > 0\):

\[\lim_{T\to\infty}\Pr\!\left(|\bar{Y}_T - \mu_Y| > e\right) \leq \frac{\sigma_Y^2}{Te^2} = 0 \qquad \square\]

Unbiasedness plus \(\text{Var}(\bar{Y}_T) = \sigma^2_Y/T \to 0\) gives
convergence in m.s., hence in probability.

{Why might Chebyshev's LLN fail for time series?}

Under independence the double-sum variance collapses:

\[\text{Var}(\bar{Y}_T) = \frac{1}{T^2}\sum_{t=1}^{T}\sum_{s=1}^{T}\text{Cov}(Y_t, Y_s) = \frac{1}{T^2}\sum_{t=1}^{T}\text{Var}(Y_t).\]

The denominator is \(O(T^2)\) and the numerator is \(O(T)\) --- so
\(\text{Var}(\bar{Y}_T) = O(1/T)\). But if we allow for an arbitrarily
high degree of temporal correlation, the cross terms no longer vanish.
The numerator is no longer \(O(T)\), and our crucial consistency result
risks falling to pieces.

Temporal dependence inflates \(\text{Var}(\bar{Y}_T)\); the numerator is
no longer \(O(T)\) and the LLN can fail.

\begin{tcolorbox}[enhanced jigsaw, arc=.35mm, bottomrule=.15mm, bottomtitle=1mm, breakable, colback=white, colbacktitle=quarto-callout-note-color!10!white, colframe=quarto-callout-note-color-frame, coltitle=black, left=2mm, leftrule=.75mm, opacityback=0, opacitybacktitle=0.6, rightrule=.15mm, title=\textcolor{quarto-callout-note-color}{\faInfo}\hspace{0.5em}{Big-O notation}, titlerule=0mm, toprule=.15mm, toptitle=1mm]

The sequence \(\{b_T\}\) is said to be \textbf{at most of order}
\(T^\alpha\), written \(b_T = O(T^\alpha)\), if there exists a finite
constant \(C > 0\) and integer \(T^*\) such that for all \(T \geq T^*\),
\(|T^{-\alpha}b_T| \leq C\).

\end{tcolorbox}

{The stark message} {The sample mean --- the most widely used statistic
in all of econometrics --- can fail to converge when observations are
temporally dependent. This is why time series analysis requires its own
toolkit.}

Dependence is not merely an inconvenience --- it strikes at the heart of
statistical inference. We will return to this when we study the ergodic
theorem.

\begin{center}\rule{0.5\linewidth}{0.5pt}\end{center}

\subsection{Review questions}\label{review-questions}

\begin{enumerate}
\def\labelenumi{\arabic{enumi}.}
\item
  Explain formally what is a random variable, a stochastic process, a
  time series, and what is meant by the analysis of time series.
\item
  Explain what is a population moment; what is the population variance
  of a random variable; and what are the covariance and correlation
  between two random variables.
\item
  State Chebyshev's LLN for an iid process. Prove it. Give an intuitive
  explanation of how and why it works. Explain its empirical relevance.
\item
  Explain with explicit reference to \(\text{Var}(\bar{Y}_T)\) why the
  LLN can potentially fail for time series data.
\end{enumerate}

\begin{tcolorbox}[enhanced jigsaw, arc=.35mm, bottomrule=.15mm, bottomtitle=1mm, breakable, colback=white, colbacktitle=quarto-callout-note-color!10!white, colframe=quarto-callout-note-color-frame, coltitle=black, left=2mm, leftrule=.75mm, opacityback=0, opacitybacktitle=0.6, rightrule=.15mm, title=\textcolor{quarto-callout-note-color}{\faInfo}\hspace{0.5em}{Note on solutions}, titlerule=0mm, toprule=.15mm, toptitle=1mm]

Separate solutions will not be provided --- answers can be found in
these notes. Use these questions to consolidate knowledge and as
discussion prompts when studying with others.

\end{tcolorbox}

\begin{center}\rule{0.5\linewidth}{0.5pt}\end{center}

\subsection{Further reading}\label{further-reading}

\begin{itemize}
\tightlist
\item
  \textbf{Statistical fundamentals} --- Appendix B and C of Wooldridge,
  \emph{Introductory Econometrics}, or any solid undergraduate
  statistics textbook.
\item
  \textbf{Main textbook} --- Shumway and Stoffer, \emph{Time Series
  Analysis and Its Applications with R Examples} (\textbf{SS}). For
  Topic 1, see SS Chapters 1.1--1.2.
\item
  The material in the stark message section is not directly from SS ---
  it represents a specific arrangement of ideas for this course.
\end{itemize}

\begin{center}\rule{0.5\linewidth}{0.5pt}\end{center}




\end{document}
